Esta interpretación la redactó Claude a partir de los registros de las corridas; no es del autor.

Luna no obtuvo \textsc{Certified} en ninguno de los 15 episodios: casi todos sus programas fueron válidos, pero ninguno bajó de $\tau_{\mathrm{ref}}$. Su mejor peor caso quedó entre un 7{,}5\,\% y un 73\,\% por encima de la referencia según la variante (el peor, norma del gradiente con $n = 6$). Sus mejores programas son sucesiones sin paso largo (el mayor de sus mejores pasos es 2{,}05). Luna sí probó pasos largos, en 10 de sus 59 programas válidos, pero como sucesiones que crecen casi geométricamente hasta pasos de 8 a 17{,}5, con peores casos cientos o miles de veces por encima de la referencia. Con las mismas instancias y el mismo juez, Astra propuso composiciones con bloques de paso $\sqrt{2}$ y un único paso moderadamente largo, entre 2{,}26 y 2{,}98, y mejoró la referencia en las cinco variantes. En este juez la diferencia entre modelos parece de conocimiento: saber que los mejores programas conocidos intercalan un paso largo, pero no demasiado, entre bloques de pasos cortos.

La memoria no cambió el resultado de Luna: le mostró su mejor programa de cada episodio anterior, que no mejoraba la referencia, y con eso no bastó. En la variante con pasos negativos permitidos, Luna no los usó en su mejor programa; su único intento con un paso fuera de rango devolvió un paso del orden de $10^{32}$ y el juez lo rechazó por inválido.
